Thus we have an explicit description of \(\operatorname{sc}(\cR)\)
in terms of the weights of \(\cR\). Likewise, one obtains a
commutative diagram
\begin{center}
\begin{tikzcd}[column sep=huge,row sep=large]
\operatorname{ad}(\cR) \arrow[rr] \arrow[d,"\wr"]
  && \cR \\
\bigl(\Gamma_0(R),\Lambda(\cR^\ast),R,R^\ast\bigr)
  \arrow[urr] &&
\end{tikzcd}
\end{center}

\begin{corollary}\label{sga3:XXI:6.6.7}
The datum \(\cR\) is simply connected if and only if
\[
 M=\Lambda(\cR).
\]
\end{corollary}

\begin{remark}\label{sga3:XXI:6.6.8}
One has
\[
 \Lambda(\cR)\cap M=\VR(R)\cap M.
\]
Consequently, \(\cR\) is semisimple if and only if
\[
 M\subset\Lambda(\cR).
\]
\end{remark}

The results of \S\ref{sga3:XXI:6.5} also give:

\begin{corollary}\label{sga3:XXI:6.6.9}
The datum \(\cR\) is the product of a simply connected datum and a
trivial datum if and only if
\[
 M\supset\Lambda(\cR).
\]
\end{corollary}

Consider now the canonical isogeny
\[
 f:\operatorname{ad}(\cR)\longrightarrow\operatorname{sc}(\cR),
\]
and put \(Z(\cR)=K(f)\). One has
\[
 Z(\cR)\simeq Z(\operatorname{sc}\cR)
 \simeq Z(\operatorname{ad}\cR).
\]

\begin{corollary}\label{sga3:XXI:6.6.10}
There is a canonical isomorphism
\[
 Z(\cR)=\Lambda(\cR)/\Gamma_0(R).
\]
More precisely, there is an exact sequence of
\(W(\cR)\)-modules
\[
 0\longrightarrow\Gamma_0(R)\longrightarrow\Lambda(\cR)
 \longrightarrow Z(\cR)\longrightarrow0.
\]
\end{corollary}

\begin{corollary}\label{sga3:XXI:6.6.11}
There is a canonical pairing
\[
 Z(\cR^\ast)\times Z(\cR)\longrightarrow\mathbb Q/\mathbb Z
\]
which puts these groups in duality.
\end{corollary}

\begin{remark}\label{sga3:XXI:6.6.12}
One has
\[
 Z(\cR\times\cR')\simeq Z(\cR)\times Z(\cR').
\]
In particular, consider simply connected data
\(\cR_i\), \(i=1,2,\ldots,n\), and a trivial datum \(\cR_0\).
Put
\[
 \cR=\cR_0\times\cR_1\times\cdots\times\cR_n.
\]
If
\[
 \cR=(M,M^\ast,R,R^\ast),
 \qquad
 \cR_0=(M_0,M_0^\ast,\varnothing,\varnothing),
\]
then
\[
 M/\Gamma_0(R)\simeq
 M_0\times Z(\cR_1)\times\cdots\times Z(\cR_n).
\]
\end{remark}

\subsection{Automorphisms}\label{sga3:XXI:6.7}

By \ref{sga3:XXI:6.1.4}, an automorphism of \(\cR\) is an
automorphism \(u\) of \(M\) such that
\[
 u(R)=R,
 \qquad
 {}^tu(R^\ast)=R^\ast.
\]
In particular, every \(w\in W(\cR)\) defines an automorphism of
\(\cR\).

\begin{lemma}\label{sga3:XXI:6.7.1}
The group \(W(\cR)\) is a normal subgroup of
\(\operatorname{Aut}(\cR)\). More precisely, if
\(u\in\operatorname{Aut}(\cR)\) and \(\alpha\in R\), then
\[
 us_\alpha u^{-1}=s_{u(\alpha)}.
\]
\end{lemma}

\begin{proof}
The proof is the same as that of \ref{sga3:XXI:1.2.10}.
\end{proof}
